\section{Methods}\label{sec:methods}

\subsection{Overview}\label{sec:overview}
For each patient $i=1,\dots,N$ and each molecular layer we form the signed change $\dchange_i(v)$ of every gene $v$ between the tumour and the matched normal specimen, place it on a fixed interaction network, and summarise its arrangement by persistent homology. The patient is the unit of inference throughout. Section~\ref{sec:ph} defines the complex, the filtrations and a tie rule; Section~\ref{sec:endpoints} the endpoints; Section~\ref{sec:inference} the tests; Section~\ref{sec:baselines} the non-topological baselines; Section~\ref{sec:simulation} the simulation study; and Section~\ref{sec:cohort} the breast cancer data.

\subsection{Network, complex and change vector}\label{sec:network}
Let $G=(V,E)$ be the STRING v12.0 human physical-interaction network \citep{szklarczyk2023} at combined score $\ge 700$, restricted to the genes measured in a layer, with the induced edges, and reduced to its largest connected component. The edges come from an external database; the vertex set is fixed by label-blind assay eligibility filters (Section~\ref{sec:cohort}), which use all specimens pooled and do not change when a patient's two specimens are exchanged. Each layer has its own network, so that no measured gene is lost to a requirement that it be measured in another layer. STRING physical links include evidence of physical proximity and shared complex membership, so an edge need not denote direct binding, and filling every 3-clique with a triangle is a modelling choice. A cycle in $K$ therefore describes the arrangement of assay values on a reference graph, not a ring of proteins in breast tissue. Let $K$ be the 2-skeleton of the clique complex of $G$: its vertices, edges and triangles (3-cliques). The 2-skeleton determines $\Hone$, so higher simplices are not needed. At this score threshold the full STRING graph has 9{,}830 genes, 85{,}576 edges and 967{,}829 triangles, and the static complex has first Betti number 4{,}160.

For patient $i$ the change vector is
\begin{equation}\label{eq:delta}
\dchange_i(v) = x_{T,i}(v) - x_{N,i}(v),
\end{equation}
where $x=\log_2(\mathrm{TPM}+1)$ for RNA and $x$ is the gene-level promoter $M$-value for methylation (Section~\ref{sec:cohort}). For RNA, $\dchange=1$ means $(\mathrm{TPM}_T+1)/(\mathrm{TPM}_N+1)=2$, not an exact two-fold TPM change. For methylation, $\dchange$ is a difference of medians over a gene's promoter probes of $M=\log_2(\beta/(1-\beta))$; for a single probe a difference of one is a doubling of the methylation odds. Exchanging the two specimens of a patient replaces $\dchange_i$ by $-\dchange_i$; this is the only property of $\dchange_i$ that the tests require. A standardised change (each specimen centred by its median across genes, then each gene scaled by the median and median absolute deviation of all $2N$ specimens, with no tumour/normal label) is a pre-specified sensitivity analysis; it is also sign-equivariant.

\subsection{Filtrations, tie rule and persistence}\label{sec:ph}
For a function $f:V\to\R$, the lower-star filtration assigns to every simplex $\sigma\in K$ the value $f(\sigma)=\max_{v\in\sigma} f(v)$ and considers the nested subcomplexes $K_t=\{\sigma: f(\sigma)\le t\}$. Faces enter no later than their cofaces, so $(K_t)$ is a filtration. We use two: the \emph{higher-change} filtration $f=-\dchange_i$, in which the most increased genes enter first, and the \emph{lower-change} filtration $f=\dchange_i$, in which the most decreased genes enter first.

\paragraph{Tie rule.} Measured changes can be tied, and persistence pairs (and cycle representatives) depend on how simplices with equal values are ordered. We rank genes by the pair $(f(v), v)$, so that ties are broken by a fixed gene index, and order simplices by the largest rank among their vertices, then by dimension, then by their sorted vertex indices. This defines one total order of $K$ without changing any value. Persistence pairs are computed in this order, and each interval's birth and death are the observed values $f(\sigma)$ of its birth edge and death triangle, so every lifetime is a difference of two observed changes. Intervals of length zero arise when a birth edge and its death triangle enter at the same vertex value. This happens through tied gene values, and also structurally, whenever the last edge of a triangle and the triangle itself share their maximal vertex; such pairs are intrinsic to lower-star filtrations. They are not intervals of the filtration by values and are discarded. Because $-\dchange$ is ranked by the same rule, the higher-change filtration of $\dchange$ coincides with the lower-change filtration of $-\dchange$, ties included.

\paragraph{Persistence.} Persistent homology over $\Fii$ \citep{edelsbrunner2002} gives, for each filtration, a diagram $D$ of $\Hone$ intervals $[b,d)$. Essential classes (cycles of the complex itself, $d=\infty$) are counted but excluded from the endpoints. This choice defines the question: the static complexes of the two layers have first Betti numbers \rnaNetBone{} and \methNetBone{}, and a network cycle that is never filled contributes nothing to the endpoints, however its genes change. We compute persistence with GUDHI \citep{maria2014} on the simplex ranks and map births and deaths back to observed values. Two standard properties are used below: for a constant $c$, adding $c$ to every value shifts every interval by $c$ and leaves every lifetime unchanged; and multiplying every value by $a>0$ multiplies every lifetime by $a$. For two functions $f,g$ on the same complex, the bottleneck distance between their diagrams is at most $\lVert f-g\rVert_\infty$ \citep{cohensteiner2007}; under the optimal matching, a matched lifetime changes by at most $2\lVert f-g\rVert_\infty$, and short intervals may be matched to the diagonal. The theorem gives no comparable guarantee for the membership of individual cycle representatives.

\paragraph{Exact cycle representatives.} For gene attribution we need, for each finite class, a cycle that represents it. We reduce, over $\Fii$ and in the total order above, the boundary columns of all triangles that kill $\Hone$ classes, including those paired at zero length. Triangles that create $\Htwo$ classes reduce to zero in the standard algorithm and never serve as pivots, so this is the standard column reduction restricted to the columns that matter for $\Hone$. The reduced column of a death triangle $\tau$ is a 1-cycle whose last edge is the birth edge of the class; it is not a boundary before $\tau$ enters and is one from $\tau$ on, so it represents exactly the class born at that edge and killed by $\tau$. Its vertices are the class's \emph{representative members}. A class has many representatives; this one is fixed by the total order and is deterministic, not biologically canonical, and the dependence of gene scores on it is measured (Section~\ref{sec:cohort}). In the software every pivot is checked against GUDHI's pairing, and unit tests check every representative against an independent full reduction of the boundary matrix.

\subsection{Endpoints}\label{sec:endpoints}
For patient $i$ let $D^{\uparrow}_i$ and $D^{\downarrow}_i$ be the finite $\Hone$ diagrams of the higher- and lower-change filtrations, and
\begin{equation}
\Lup_i=\sum_{[b,d)\in D^{\uparrow}_i}(d-b),\qquad \Ldown_i=\sum_{[b,d)\in D^{\downarrow}_i}(d-b).
\end{equation}
The \emph{primary endpoint} is the asymmetry $A_i=\Lup_i-\Ldown_i$.

\begin{proposition}[Exact sign reversal]\label{prop:antisym}
For every change vector $\dchange$, $A(-\dchange)=-A(\dchange)$.
\end{proposition}
\begin{proof}
The higher-change filtration of $-\dchange$ uses $f=\dchange$, ranked by $(\dchange(v),v)$; this is the lower-change filtration of $\dchange$, with the same total order and the same observed values. Hence $D^{\uparrow}(-\dchange)=D^{\downarrow}(\dchange)$ and, symmetrically, $D^{\downarrow}(-\dchange)=D^{\uparrow}(\dchange)$, so $\Lup$ and $\Ldown$ exchange and $A$ changes sign.
\end{proof}

A positive $A_i$ means a greater total finite $\Hone$ lifetime in the higher-change filtration than in the lower-change filtration; it need not mean more cycles. Two variants are pre-specified. The \emph{rank variant} replaces $\dchange_i$ by its centred within-patient average rank $\rho_i(v)=\rank_i(v)-(|V|+1)/2$; average ranks satisfy $\rank(-\dchange)=|V|+1-\rank(\dchange)$ exactly, so $\rho$ is also exactly sign-reversing, and a result that holds for $\dchange$ but not for $\rho$ is driven by the size of the changes rather than their order. The \emph{direction-free total} $P_i=\Lup_i+\Ldown_i$ is invariant under exchange and is used only descriptively.

\paragraph{Gene score.} With $C^{\uparrow}_i$ and $C^{\downarrow}_i$ the finite classes of the two filtrations and $z_c$ the representative of class $c$,
\begin{equation}
a_i(g)=\sum_{c\in C^{\uparrow}_i:\,g\in\supp z_c}\pers(c)\;-\sum_{c\in C^{\downarrow}_i:\,g\in\supp z_c}\pers(c),\qquad a_i(g;-\dchange_i)=-a_i(g;\dchange_i).
\end{equation}
A gene with a large score is a \emph{topological participant}: it lies on persistent cycles of one direction. It need not itself be strongly changed. Each class contributes its whole lifetime to every vertex of its representative, so $a_i(g)$ is a participation measure and not an additive decomposition of $A_i$: summing $a_i(g)$ over genes weights each class by the size of its representative.

\paragraph{Network-adjusted score.} To separate the arrangement of changes on the network from their size, we compare each patient's asymmetry with a reference in which that patient's own changes are placed at other positions of the same kind. Genes are stratified by degree decile $\times$ local clustering coefficient tercile, optionally crossed with terciles of each gene's mean and variance over all $2N$ specimens pooled (pooled statistics do not change when a patient's specimens are exchanged). For patient $i$, $R$ permutations $\pi_{i1},\dots,\pi_{iR}$ of the genes within strata are drawn from a random seed that depends only on a master seed and on $i$, never on the data. Relabelling $r$ moves $\dchange_i(g)$ to node $\pi_{ir}(g)$, and $A$ is recomputed, giving $A^{*}_{ir}$ with mean $\mu_i$ and standard deviation $\sigma_i$. The network-adjusted score is
\begin{equation}\label{eq:Z}
Z_i=\frac{A_i-\mu_i}{\max(\sigma_i,\varepsilon)},\qquad \varepsilon=0.1\times\operatorname{median}_j\sigma_j .
\end{equation}

\begin{proposition}\label{prop:Z}
For fixed strata, $Z_i(-\dchange_i)=-Z_i(\dchange_i)$ for every patient.
\end{proposition}
\begin{proof}
The permutations do not depend on $\dchange_i$, and relabelling commutes with negation, so by Proposition~\ref{prop:antisym} $A^{*}_{ir}(-\dchange_i)=-A^{*}_{ir}(\dchange_i)$ for every $r$. Hence $\mu_i$ changes sign, $\sigma_i$ is unchanged, $\varepsilon$ (a function of the $\sigma_j$) is unchanged, and $Z_i$ changes sign.
\end{proof}

The floor $\varepsilon$ prevents a near-zero reference spread from inflating $Z_i$; if every $\sigma_j$ is zero, the floor is set to the smallest positive double-precision number, so $Z_i$ is defined (and equals zero whenever $A_i=\mu_i$). In the studentized statistics below, $t=0$ when the sample standard deviation is zero. What $Z_i$ removes is determined by the reference: it adjusts for the distribution of the patient's own changes and for degree and clustering, but not for the adjacency of the changed genes to one another. This limitation matters in practice (Section~\ref{sec:res-hub}).

\subsection{Statistical inference}\label{sec:inference}
\paragraph{Level 1: paired asymmetry.} The null hypothesis $H_{\mathrm{sym}}$ states that, within each patient, the tumour and normal profiles are exchangeable, and that patients are independent. Under $H_{\mathrm{sym}}$ the vector $(A_1,\dots,A_N)$ has the same distribution as $(s_1A_1,\dots,s_NA_N)$ for every fixed sign vector $s\in\{-1,1\}^N$ (Proposition~\ref{prop:antisym}). For each layer we use the studentized mean $t=\sqrt{N}\,\bar A/s_A$ and the Monte Carlo sign-flip $p$-value
\begin{equation}
p=\frac{1+\#\{b\le B:|t_b|\ge|t_0|\}}{B+1},\qquad B=9{,}999,
\end{equation}
with one sign matrix shared by all layers. Counting the observed statistic makes $P(p\le\alpha)\le\alpha$ for every $\alpha$ under $H_{\mathrm{sym}}$ when the sign vectors are drawn at random, so the test is valid in finite samples without enumerating all $2^{N}$ sign vectors \citep{phipson2010}. Holm's procedure \citep{holm1979} over layers controls the family-wise error rate (FWER) at 0.05. $H_{\mathrm{sym}}$ can fail for reasons that are not biology of interest, for example if tumour and normal specimens were processed in different batches; this is examined in sensitivity analyses.

\paragraph{Level 2: network-adjusted asymmetry.} $Z_1,\dots,Z_N$ are tested in the same way, with Holm over layers. By Proposition~\ref{prop:Z} the test is valid under joint sign symmetry of $(Z_1,\dots,Z_N)$, which holds under $H_{\mathrm{sym}}$ because the strata use pooled statistics and the floor $\varepsilon$ is invariant under sign changes; the shared floor makes the $Z_i$ dependent, which joint sign symmetry allows. Whether a rejection reflects arrangement on the network, rather than an effect that the reference does not remove, is a property of the reference model and not of the test; Section~\ref{sec:res-hub} checks it.

\paragraph{Level 3: genes.} The family is every (layer, gene) with $a_i(g)\ne0$ in at least 10 patients. Each gene receives a two-sided sign-flip $p$-value of its studentized mean ($2\times10^{6}$ flips in the real data, $2\times10^{5}$ in simulations, so that the smallest attainable $p$-value lies below the first Holm threshold), and Holm's procedure is applied over the whole family. The selection depends only on the $|a_i(g)|$, which joint sign flips of all genes do not change. Under $H_{\mathrm{sym}}$ for every patient (the complete null), the $p$-values are therefore valid conditionally on the family, and Holm controls the FWER. Under a partial null, two further points would need justification. A gene's score depends on its neighbours through the cycle representatives, so a gene whose own change is symmetric can have an asymmetric score when its neighbours change; the null hypothesis that is tested is sign symmetry of the score $a_i(g)$, not of the gene's own change. And we have not shown that the $p$-value of a gene with a sign-symmetric score remains valid conditionally on a family selected using all genes. We therefore claim family-wise error control for the gene family under the complete null, which the simulations check (Table~\ref{tab:calibration}), and report the genes as Holm-adjusted participants rather than as discoveries with strong FWER control.

\subsection{Baselines}\label{sec:baselines}
Every baseline is computed on the same $\dchange_i$, network, patients and sign flips. With $\dchange^{+}=\max(\dchange,0)$, $\dchange^{-}=\max(-\dchange,0)$, a threshold $\tau=1$ and $N(v)$ the neighbours of $v$:
\begin{itemize}
\item[B1] signed burden $\#\{v:\dchange(v)\ge\tau\}-\#\{v:\dchange(v)\le-\tau\}$;
\item[B2] signed autocorrelation $I(\dchange^{+})-I(\dchange^{-})$, with Moran's $I$ \citep{moran1950} on the binary adjacency of $G$ and $I=0$ for a variable with zero variance;
\item[B3] the size of the largest connected component of $\{\dchange\ge\tau\}$ minus that of $\{\dchange\le-\tau\}$;
\item[B4] the $\Hzero$ asymmetry, total finite $\Hzero$ persistence of the higher- minus the lower-change filtration;
\item[B5] Moran's $I(\dchange)$, direction-free and descriptive;
\item[B6] single-gene paired tests: the sign-flip test of the studentized mean of $\dchange_i(g)$ for every gene of the layer network, with Holm over all network genes (a comparison restricted to the topology family is reported as well);
\item[B7] neighbourhood-smoothed single-gene tests: the same test applied to the mean of $\dchange_i$ over $g$ and its neighbours, with Holm over all network genes;
\item[B8] specimen-level $\Hone$: total finite $\Hone$ persistence of each specimen's own standardised values in the superlevel filtration (highest values first), tumour minus normal;
\item[B9] local-extremum depth $\sum_v\max(0,\min_{u\in N(v)}\dchange(u)-\dchange(v))-\sum_v\max(0,\dchange(v)-\max_{u\in N(v)}\dchange(u))$.
\end{itemize}
B9 asks, without topology, whether some genes lie below all their changed neighbours, which is the configuration at the centre of a higher-change cycle. The directional baselines B1--B4 and B9 reverse sign exactly under exchange and enter the same sign-flip test.

\subsection{Simulation study}\label{sec:simulation}
All simulated data are synthetic. Networks are breadth-first balls of 800 genes (400 for the costly network-adjusted runs) around random seed genes of the real STRING graph at score $\ge700$. Every edge of a ball is a real STRING edge, and the balls have the high clustering and triangle density of the real network, but a truncated ball need not reproduce the full graph's degree or triangle distribution: genes on its boundary lose neighbours. For patient $i$ and gene $v$,
\begin{equation}
x_{T,i}(v)=\sqrt{s}\,U_{i}(v)+\sqrt{1-s}\,E_{T,i}(v),\qquad x_{N,i}(v)=\sqrt{s}\,U_{i}(v)+\sqrt{1-s}\,E_{N,i}(v),
\end{equation}
with a shared patient component $U$, $s=0.2$, $N=78$ patients, and independent errors drawn from a standard normal, a Student $t_3$ scaled to unit variance, or a centred, unit-variance gamma(2) distribution, or normal errors in which 5\% of specimens have five times the spread (outliers). Under the null the two specimens of each patient are exchangeable for every noise type. Two features of this model limit what it tests. The shared component $U_i$ cancels exactly from the plain change \eqref{eq:delta}, and the difference of two independent, identically distributed gamma errors is symmetric, so with the plain change the skewed noise does not produce a skewed change; the skewed and $t_3$ runs used the standardised change, in which pooled centring and scaling enter. The model has no correlation between genes and no patient-specific heterogeneity of the change. Each run uses its own pre-assigned range of seeds. Development runs were used while the method was being finalised; the main results were then repeated once on a reserved verification range with the final code, and the verification results are the ones reported in the main text.

\paragraph{Calibration.} With no change at all, we record whether Level 1 and the gene family reject at 0.05 (single layer). A run fails the inflation gate if the one-sided exact binomial test of ``rejection rate $\le0.05$'' gives $p<0.01$. No complete-null run of the Level 2 test with the final reference was made; its validity rests on Proposition~\ref{prop:Z} and on the implementation checks shared with Level 1. Verification runs used the plain change \eqref{eq:delta} with normal and outlier noise; development runs used the standardised change with normal, $t_3$ and skewed noise. Sign-flip validity does not depend on which sign-equivariant change is used.

\paragraph{Matched cycle--module experiment.} On each 400-gene ball we choose up to three centres $c$ of degree 4--20 whose neighbours alone enclose a hole (the clique complex of the subgraph induced by $N(c)$ has first Betti number at least one), with closed neighbourhoods that neither overlap nor are joined by an edge. A ball with no such centre is discarded and a new ball drawn; the number of discarded balls was not recorded. From a single draw of noise we build two cohorts. In the \emph{cycle} arm every neighbour of every centre is increased by a fixed size in the tumour, and the centre is not. In the \emph{module} arm the centre and all but one randomly chosen neighbour are increased by the same size. The two arms have the same network region, the same number of changed genes, the same change sizes and the same noise. They differ in which genes are changed: the centre and one neighbour exchange roles. This exchange leaves the changed genes enclosing an unchanged gene in the cycle arm, but it also changes the degrees and clustering coefficients of the changed genes, the edges and triangles among them, how they attach to the rest of the graph, and their local autocorrelation. The experiment therefore measures the effect of exchanging a changed neighbour for the centre in a matched network region, not the effect of a hole alone.

\paragraph{What the two arms do to the planted subgraph.} In the noiseless planted subgraph the two arms differ topologically. In the cycle arm the changed genes are $N(c)$, whose induced clique complex has $\beta_1\ge1$ by the selection rule. In the module arm the changed genes are $S\cup\{c\}$ with $S\subset N(c)$, and $c$ is adjacent to every vertex of $S$, so every clique of $S$ extends to a clique containing $c$; the induced clique complex is the cone over the clique complex of $S$ with apex $c$, which is contractible, so its $\beta_1$ is zero. This statement concerns the induced planted subgraph only. Our endpoint is total finite persistence on the whole noisy network, where a planted cycle can be filled through vertices outside the planted set (for example, an outside vertex adjacent to every vertex of the rim fills it before the centre enters), where the selection rule allows more than one planted class per centre, and where many other classes change between the arms. We did not measure how often a planted class survives in the full filtration. For $A_i$, its rank variant, $Z_i$ ($R=19$), B1--B4 and B9 we record the studentized statistic $t$ in both arms. The pre-declared summary is the area under the curve across data sets, $\mathrm{AUC}=P(t_{\mathrm{cycle}}>t_{\mathrm{module}})+\tfrac12P(t_{\mathrm{cycle}}=t_{\mathrm{module}})$, computed over all pairs of a cycle-arm and a module-arm data set (the Mann--Whitney statistic divided by $n^2$); leaving out each data set's own matched counterpart changes it by at most 0.002. The plan stated that $\Hone$ adds information here if $A_i$ or $Z_i$ separates the arms better than B2, B3, B4 and B9. An AUC below 0.5 means that the statistic is larger in the module arm, and its \emph{oriented} AUC is $\max(\mathrm{AUC},1-\mathrm{AUC})$. As secondary summaries we report, within matched pairs, the share of data sets with $t_{\mathrm{cycle}}>t_{\mathrm{module}}$ (wins), the number of ties, and the mean paired difference $\Delta t=t_{\mathrm{cycle}}-t_{\mathrm{module}}$ with a normal-approximation 95\% interval. The two summaries answer different questions: the paired share asks how a statistic responds to the exchange when the network and the noise are held fixed, and the AUC asks whether it orders data sets drawn from different networks. Sizes 1 and 2 (in units of the noise standard deviation) were fixed in advance, with 200 data sets each. The plain change \eqref{eq:delta} was used.

\paragraph{Change at hubs (N3).} On 400-gene balls, the genes in the top 5\% of degree are increased by 1 in every tumour, so that topology changes only as a consequence of strong change at high-degree genes. We record the rejection rate of the Level 2 test of $Z_i$ ($R=19$, strata by degree decile $\times$ clustering tercile; 500 data sets). This is not a calibration run: the change is real. A rate near 0.05 means that the reference absorbs this structured alternative, and a higher rate means that $Z_i$ would attribute it to arrangement.

\subsection{Breast cancer data}\label{sec:cohort}
The development cohort is TCGA-BRCA \citep{tcga2012}, open-access files of the NCI Genomic Data Commons: RNA-seq (STAR, TPM) and DNA methylation (Illumina 450K, SeSAMe $\beta$ values). We include female patients with one primary tumour and one solid normal tissue specimen in both layers, excluding FFPE specimens; if a patient has several eligible files of one type, the file whose aliquot barcode sorts first is used. The GDC was queried on 16 November 2026 (query log and file manifest in \texttt{docs/provenance}): 1{,}097 patients had RNA or methylation files of these types, 1{,}084 of them female, and \cohortN{} met all criteria. All 437 GDC annotations of these patients \citep{gdcannotations} were reviewed before any molecular value was opened; none concerns an RNA or methylation file or aliquot used here (they concern somatic-mutation calls and exome DNA aliquots), so no patient was excluded. All preprocessing is label-blind: every filter uses the $2N$ specimens pooled. RNA: protein-coding genes (GENCODE gene type) with TPM $\ge1$ in at least half of all specimens, $x=\log_2(\mathrm{TPM}+1)$; \rnaGenes{} genes. Gene symbols were matched to the preferred names of STRING proteins; duplicate symbol pairs were collapsed into one edge, and genes absent from STRING or outside the largest connected component were left out of the network. Methylation: probes in the general mask of \citet{zhou2017}, probes on chromosomes X, Y and M or unplaced, probes absent from the EPIC array (so that the confirmation cohort can use the same probes) and probes missing in more than 10\% of specimens are removed; a probe is in the promoter of a gene if a protein-coding transcript of that gene starts 0--1500 bp downstream of the CpG on the transcript's strand (the TSS200 and TSS1500 regions; the annotation's signed distance to the transcription start site is strand-aware), using the GENCODE v41 probe annotation of the same authors, and a probe in the promoters of several genes counts for each; $\beta$ is clipped to $[0.001,0.999]$, $M=\log_2(\beta/(1-\beta))$, and a gene's value is the median $M$ over its promoter probes (at least two; \methPromoterProbes{} probes in all), with the probe list per gene frozen and hashed. Genes with no value in some specimen are dropped, leaving \methGenes{} genes; nothing is imputed. The layer networks have \rnaNetV{} genes, \rnaNetE{} edges and \rnaNetT{} triangles (RNA; first Betti number \rnaNetBone{}) and \methNetV{} genes, \methNetE{} edges and \methNetT{} triangles (methylation; \methNetBone{}). Pre-specified sensitivity analyses repeat Level 1 (Level 2 is computed for the primary analysis only, because each $Z_i$ analysis needs about 15{,}000 persistence computations per layer) with the rank variant, with the standardised change, with STRING score thresholds 400 and 900, with composition adjustment, and within pairs whose two specimens share a plate; the gene scores are recomputed with a reversed tie rule. For composition adjustment, each specimen's stromal and immune proxies are the mean standardised RNA value over the ESTIMATE stromal and immune gene sets \citep{yoshihara2013} (\estimateStromal{} and \estimateImmune{} of the 141 genes of each set pass the RNA filters; these are proxies, not ESTIMATE scores). The tumour-minus-normal differences $c_i\in\R^2$ are regressed out of $\dchange_i$ gene by gene, through the origin, and the same RNA-derived proxies are used for the methylation layer. The adjustment keeps the sign reversal: with $X$ the $N\times2$ matrix of the $c_i$, $Y$ the $N\times|V|$ matrix of the $\dchange_i$ and $S$ a diagonal matrix of patient signs, the fit to $(SX,SY)$ has the same coefficients $\hat B=(X^\top X)^{-1}X^\top Y$ (because $S^\top S=I$), so the residuals are $S(Y-X\hat B)$. Any covariate that reverses sign with the patient's specimens can be included in this way; a covariate that does not would break the argument. The proxy gene sets overlap the analysed RNA genes, so the adjustment can remove part of a signal carried by those genes. An independent cohort is planned for confirmation (Section~\ref{sec:limitations}). The full plan is in Supplement~\ref{sec:plan}. Two exploratory analyses were added after the frozen results had been seen: the Level~2 test of $Z_i$ after the same composition adjustment (\texttt{scripts/tcga\_exploratory.py}), and a count of edges among the highest-degree genes of the STRING graph compared with degree-preserving rewiring (\texttt{scripts/hub\_check.py}, no patient data). Supplementary Table~\ref{tab:prespec} lists every analysis by endpoint as pre-specified or exploratory.
